\emph{Preliminaries.} By Fact \claimref{field-theory.imported-tower-limit} the null pairs of distinct planes anticommute and $\{v_{j}, w_{j}\} = 1$, $v_{j}^{2} = w_{j}^{2} = 0$; with $v_{j} = \tfrac12(e_{j} + f_{j})$, $w_{j} = \tfrac12(e_{j} - f_{j})$ this is the product in $\Cl(m,m)$, and $n_{j} = v_{j}w_{j} = \tfrac12(1 - e_{j}f_{j})$ is an idempotent with $w_{j}v_{j} = 1 - n_{j}$. Two facts are used throughout. (a) Elements of even degree built from the generators of disjoint sets of planes commute, since each generator of one anticommutes with each generator of the other and an even number of exchanges is made. (b) For such elements $x, y$ the scalar part factorises, $\langle xy\rangle_{0} = \langle x\rangle_{0}\langle y\rangle_{0}$, since a product of blades on disjoint index sets is a blade, scalar only when both are. Hence for $\rho = \prod_{j}(p_{j}n_{j} + (1 - p_{j})(1 - n_{j}))$ and $\langle n_{j}\rangle_{0} = \tfrac12$, $\omega(1) = 1$, $\omega(n_{j}) = p_{j}$, and $\omega(X_{1}\cdots X_{r}) = \prod\omega(X_{i})$ for even $X_{i}$ on disjoint sets of planes. The scalar part is cyclic, $\langle xy\rangle_{0} = \langle yx\rangle_{0}$, so $\omega([X, Y]) = 2^{m}\langle[\rho, X]Y\rangle_{0}$; and $\omega(X^{*}) = \omega(X)$ for the self-adjoint $\rho$ of a real state, since $\langle x^{*}\rangle_{0} = \langle x\rangle_{0}$.

\emph{(i)} For $i \ne j$, $\{w_{i}, w_{j}\} = 0$ and $w_{i}^{2} = 0$, so $S^{2} = N^{-1}\big(\sum_{i}w_{i}^{2} + \sum_{i<j}\{w_{i}, w_{j}\}\big) = 0$, and $\{S, S^{*}\} = N^{-1}\sum_{i,j}\{w_{i}, v_{j}\} = N^{-1}\sum_{i}1 = 1$. Hopping and pairing elements and their adjoints are of degree two in the generators of their own block, so (a) gives the commutation across blocks. The check verifies $S^{2} = 0$ and $\{\sqrt{N}S, \sqrt{N}S^{*}\} = N$ for $N = 1, \dots, 4$ and the commutation for two and three blocks.

\emph{(ii)} Moving $w_{1}$ to the left past $v_{2}$ and $w_{2}$ (two exchanges), $bb^{*} = v_{1}w_{2}v_{2}w_{1} = v_{1}w_{1}w_{2}v_{2} = n_{1}(1 - n_{2})$, and likewise $b^{*}b = v_{2}w_{1}v_{1}w_{2} = n_{2}(1 - n_{1})$, so $[b, b^{*}] = n_{1} - n_{2}$. By (a) the cross terms of $[B, B^{*}]$ vanish and $[B, B^{*}] = N^{-1}\sum_{i}(n_{1i} - n_{2i}) = -2K_{3}/N$. Since $[n_{j}, v_{j}] = v_{j}$ and $[n_{j}, w_{j}] = -w_{j}$, $[K_{3}, b_{i}^{*}] = b_{i}^{*}$ and $[K_{3}, b_{i}] = -b_{i}$: the triple closes as $\mathfrak{su}(2)$. $\Omega_{\mathrm{pol}}$ is annihilated by every $b_{i}$ (each upper mode is empty) and has $K_{3} = -N/2$, so it is a lowest-weight vector of spin $N/2$; $(B^{*})^{k}\Omega_{\mathrm{pol}}$ has $K_{3} = -N/2 + k$, on which $[B, B^{*}] = 1 - 2k/N$. Writing $\psi_{k} = (\sqrt{N}B^{*})^{k}\Omega_{\mathrm{pol}}$, $\sqrt{N}B\,\psi_{k+1} = \big(\sqrt{N}B^{*}\sqrt{N}B + N[B, B^{*}]\big)\psi_{k}$ gives by induction $\sqrt{N}B\,\psi_{k+1} = (k + 1)(N - k)\psi_{k}$, nonzero for $k < N$, so $\psi_{0}, \dots, \psi_{N}$ are nonzero; $\psi_{N+1}$ would have $K_{3} = N/2 + 1$, above the largest eigenvalue $N/2$ of $K_{3}$, so it is $0$. The canonical relation is the limit $1 - 2k/N \to 1$ at fixed $k$. The reflection $f_{1i} \mapsto -f_{1i}$ preserves $Q$, so it extends to an automorphism of the Clifford algebra; it exchanges $v_{1i}$ and $w_{1i}$, maps $b_{i} = v_{1i}w_{2i}$ to $w_{1i}w_{2i}$ and $n_{1i}$ to $1 - n_{1i}$, hence $N[B, B^{*}] = \sum_{i}(n_{1i} - n_{2i})$ to $N - \sum_{i}(n_{1i} + n_{2i})$, the commutator of the pairing sum of Proposition \claimref{field-theory.how-close-tower-comes-ccr}. The check verifies the commutator, the $\mathfrak{su}(2)$ relations, the image under the reflection and the ladder for $N = 1, \dots, 4$.

\emph{(iii)} As in (ii), $cc^{*} = w_{1}w_{2}v_{2}v_{1} = w_{1}(1 - n_{2})v_{1} = (1 - n_{1})(1 - n_{2})$ and $c^{*}c = v_{2}v_{1}w_{1}w_{2} = n_{1}n_{2}$, so $[c, c^{*}] = 1 - n_{1} - n_{2}$. By the preliminaries, $\omega([b, b^{*}]) = p_{1} - p_{2}$ and $\omega([c, c^{*}]) = 1 - p_{1} - p_{2}$. For the Gibbs state, $n^{2} = n$ gives $e^{xn} = 1 + (e^{x} - 1)n$, so $e^{-\varepsilon(n_{2} - n_{1})/T} = (1 + (y - 1)n_{1})(1 + (y^{-1} - 1)n_{2})$, a product of the form above after normalisation, with $p_{1} = y/(1 + y)$ and $p_{2} = y^{-1}/(1 + y^{-1}) = 1/(1 + y)$. Then $p_{1} - p_{2} = (y - 1)/(y + 1) = \tanh(\varepsilon/2T)$ and $p_{1} + p_{2} = 1$; at $y = 3$ the value is $\tfrac12$. The check computes these values in $36$ product states and in the Gibbs states at $y = 2, 3, 5$, building the Gibbs element from $e^{xn} = 1 + (y - 1)n$.

\emph{(iv)} Let $\ket{s}$, $s \in \{00, 10, 01, 11\}$, be the occupation basis, $A_{s} = v_{1}^{s_{1}}v_{2}^{s_{2}}$, $P_{0} = w_{1}v_{1}w_{2}v_{2}$ (the projector on the empty state), and $E_{st} = A_{s}P_{0}A_{t}^{*}$. These satisfy $E_{st}E_{t'u} = \delta_{tt'}E_{su}$ (verified in the check for all $256$ products), so they are matrix units of $\Mat_{4}(\R)$, and the eight with $s, t$ of equal number parity span $\mathcal{A}_{e} = \Mat_{2}(\R)\oplus\Mat_{2}(\R)$. In each summand the commutators span the matrices of trace zero ($E_{st} = [E_{st}, E_{tt}]$ for $s \ne t$, $E_{ss} - E_{tt} = [E_{st}, E_{ts}]$), six dimensions in all (computed). So $\sigma_{\omega} \equiv 0$ if and only if $\omega(E_{st}) = 0$ for $s \ne t$ and $\omega$ is constant on the diagonal units of each summand, which is $\omega|_{\mathcal{A}_{e}} = \lambda\tau_{+} + (1 - \lambda)\tau_{-}$. A pure state of the block is a vector state $\psi = \sum_{s}\psi_{s}\ket{s}$ ($\Mat_{4}(\R)$ is simple); tracial would require $\psi_{00}^{2} = \psi_{11}^{2} = \lambda/2$, $\psi_{10}^{2} = \psi_{01}^{2} = (1 - \lambda)/2$ and $\psi_{00}\psi_{11} = \psi_{10}\psi_{01} = 0$, so $\lambda = 0$ and $\lambda = 1$ at once, impossible. For a product state $\rho$ is diagonal, $\rho = \mathrm{diag}(a, \beta, \gamma, d)$ with $a = (1 - p_{1})(1 - p_{2})$, $\beta = p_{1}(1 - p_{2})$, $\gamma = (1 - p_{1})p_{2}$, $d = p_{1}p_{2}$, and lies in $\mathcal{A}_{e}$. By cyclicity $\sigma_{\omega}(X, Y) = 2^{m}\langle[\rho, X]Y\rangle_{0}$, and the trace pairing is nondegenerate on each $\Mat_{2}(\R)$, so the radical of $\sigma_{\omega}$ is the commutant of $\rho$ in $\mathcal{A}_{e}$: in the even-number summand it is the diagonal ($2$ dimensions) when $a \ne d$ and everything when $a = d$, and likewise in the odd-number summand with $\beta, \gamma$. Hence the rank is $2[a \ne d] + 2[\beta \ne \gamma]$, and $a - d = 1 - p_{1} - p_{2}$, $\beta - \gamma = p_{1} - p_{2}$. Both vanish only at $p_{1} = p_{2} = \tfrac12$; for $p_{1}, p_{2} \in \{0, 1\}$ exactly one of them is nonzero. The check computes the $8 \times 8$ matrix $\omega([E_{a}, E_{b}])$ and its rank for $26$ product states, the tracial states at $\lambda = 0, \tfrac13, 1$ and the pure state $\tfrac{1}{\sqrt2}(\ket{00} + \ket{11})$ (rank $2$).

\emph{(v)} From $[v_{a}w_{b}, v_{c}w_{d}] = \delta_{bc}v_{a}w_{d} - \delta_{ad}v_{c}w_{b}$ (the CAR), with $\rho_{q}^{*} = \sum_{l}v_{l+q}w_{l}$,
\[
[\rho_{q}, \rho_{q}^{*}] = \sum_{k=0}^{M-1-q}n_{k} - \sum_{k=q}^{M-1}n_{k} = \sum_{k<q}n_{k} - \sum_{k \ge M-q}n_{k},
\]
the form stored in Theorem \claimref{hawking.current-commutator-exactly-edge-split} and Theorem \claimref{synthesis.floor-regulator-keeps-the-charge-and-the-central-term-below-the-cut-off}(ii). It is diagonal in the occupation basis, with eigenvalue $d(s) = \sum_{k<q}s_{k} - \sum_{k \ge M-q}s_{k}$ on $\ket{s}$; for $2q \le M$ the two edge sets are disjoint, so $d(s) \le q$ with equality exactly on edge-clean $s$, and since the occupation basis is a basis of eigenvectors, the operator acts on $\psi$ as $q$ exactly when $\psi$ lies in the span of the edge-clean $\ket{s}$. The product-state value follows from $\omega(n_{k}) = p_{k}$; with $p_{k} = c$ for all $k$ the two sums are $qc$ each, while the block form of two modes at occupation $c$ has rank $2[2c \ne 1]$ by (iv). The check verifies the operator identity for $M = 4, 6$ and all $q$, the eigenvalue and the converse on all $2^{M}$ basis vectors, and the values on the filled sea, at constant occupations $\tfrac12$ and $\tfrac14$, and for the Fermi-Dirac profile $p_{k} = (1 + 2^{2k - M + 1})^{-1}$.

\emph{(vi)} \emph{Restriction.} By Theorem \claimref{synthesis.null-vacuum-on-a-horizon-generator}(ii), $\omega_{0}(a^{*}(f)a(g)) = \ip{g}{Ef}$ with $E = \chi_{(-\infty,0)}(P)\otimes\id_{K}$. The $g_{k}$ are orthonormal (disjoint spectral supports), and $Eg_{k} = g_{k}$ for $k < 0$ (support in $(-\infty, 0)$) and $Eg_{k} = 0$ for $k \ge 0$ (support in $[0, \infty)$); so the two-point matrix of $\omega_{0}$ on the $g_{k}$ is diagonal with entries $1$ for $k < 0$ and $0$ for $k \ge 0$. For a symbol $\phi(P)\otimes\id_{K}$, acting in the spectral representation as multiplication by $\phi$, $\ip{g_{l}}{\phi(P)g_{k}} = \delta_{kl}\Delta^{-1}\int_{k\Delta}^{(k+1)\Delta}\phi$. A gauge-invariant quasi-free state is the state whose expectations of monomials $a^{*}(f_{1})\cdots a^{*}(f_{r})a(g_{r'})\cdots a(g_{1})$ vanish for $r \ne r'$ and are the determinants $\det[\omega(a^{*}(f_{i})a(g_{j}))]$ for $r = r'$; these depend only on the two-point function on the vectors involved, so on the CAR algebra of finitely many $g_{k}$ the restriction is fixed by the diagonal matrix $\mathrm{diag}(p_{k})$. The product state with occupations $p_{k}$ has exactly these expectations: monomials with $r \ne r'$ change the occupation numbers and have expectation $0$, and by (b) a balanced monomial factorises over modes into factors $\omega(n_{k}) = p_{k}$ with the sign of the permutation, giving the determinant of $\mathrm{diag}(p_{k})$ restricted to the index sets (verified in the check for all monomials of degree up to $3 + 3$ on four modes with occupations $\tfrac13, 1, 0, \tfrac34$). Monomials span the window's algebra, so the two states agree: the restriction is the product state, and for $\omega_{0}$ it is the filled sea of the window.

\emph{Window forms.} By (v) at $p_{k} = 1$ ($k < 0$), $0$ ($k \ge 0$), on the window of $2N$ modes $\omega_{0}([\rho_{q}, \rho_{q}^{*}]) = \min(q, 2N - q) > 0$. For $q \ne q'$ the commutators $[\rho_{q}, \rho_{q'}^{*}]$ and $[\rho_{q}, \rho_{q'}]$ are, by the same CAR identity, sums of bilinears $v_{a}w_{b}$ with $a \ne b$ (the latter is in fact $0$: the two index ranges coincide), whose expectation in a product state is $0$ by the determinant formula. So on the basis $\rho_{1}, \rho_{1}^{*}, \dots, \rho_{2N-1}, \rho_{2N-1}^{*}$ the matrix of $\sigma_{\omega_{0}}$ is block diagonal with $2 \times 2$ blocks of entries $\pm\min(q, 2N - q)$, so it is nondegenerate. A block of two window modes carries a pure product state, of rank $2$ by (iv); with symbol $\tfrac12$ every $p_{k} = \tfrac12$ and both forms vanish by (iv) and (v). The check computes the band averages of $\chi_{(-\infty, 0)}$ over $[k, k + 1)$, the $(4N - 2) \times (4N - 2)$ density form and its rank, and the block forms, for $N = 1, 2, 3$, and at occupations $\tfrac12$ and $\tfrac14$.

\emph{The infinite band.} The restriction of $\omega_{0}$ to the CAR algebra of all $g_{k}$, $k \in \mathbb{Z}$, has the same expectations on every window, so it is the product state with $n_{k} = 1$ for $k < 0$ and $0$ for $k \ge 0$. In its GNS representation with cyclic vector $\Omega$, $\|w_{k}\Omega\|^{2} = \omega_{0}(v_{k}w_{k}) = 0$ for $k \ge 0$ and $\|v_{k}\Omega\|^{2} = \omega_{0}(w_{k}v_{k}) = 1 - \omega_{0}(n_{k}) = 0$ for $k < 0$, so $w_{k}\Omega = 0$ for $k \ge 0$ and $v_{k}\Omega = 0$ for $k < 0$. Let $\psi = x\Omega$ with $x$ a polynomial in the modes $|k| < L$. In $\rho_{q}\psi = \sum_{k}v_{k}w_{k+q}\psi$ every term with $k < -L$ or $k + q \ge L$ vanishes: if $k + q \ge L$, $w_{k+q}$ anticommutes past $x$ (degree by degree) and kills $\Omega$ (the mode is empty); if $k < -L$, then $v_{k}$ anticommutes past $w_{k+q}$ ($k \ne k + q$) and past $x$ and kills $\Omega$ (the mode $k$ is filled). So $\rho_{q}$ is a finite sum on $\psi$, equal to the window operator of any window containing $[-L, L + q)$, and $\rho_{q}\psi$, $\rho_{q}^{*}\psi$ lie again in $\mathcal{D}$; in particular $\rho_{q}\Omega = 0$, the lowering operators annihilating the filled sea of every window. Take a window $[-L', L')$ with $L' \ge L + q + q'$. Then $[\rho_{q}, \rho_{q'}^{*}]\psi$ is the window commutator applied to $\psi$, and $\psi$ is edge-clean at depth $\max(q, q')$ for that window. For $q = q'$, (v) gives $[\rho_{q}, \rho_{q}^{*}]\psi = q\psi$. For $q \ne q'$ the window commutator is a difference of bilinears $v_{a}w_{b}$, $a \ne b$, with $a, b$ both among the lowest or both among the highest $\max(q, q')$ window modes (Theorem \claimref{hawking.current-commutator-exactly-edge-split}; it also follows from the CAR identity above by comparing the index ranges); on an edge-clean $\psi$ each such bilinear gives $0$ ($v_{a}w_{b}\psi = -w_{b}v_{a}\psi = 0$ for filled $a$; $w_{b}\psi = 0$ for empty $b$). And $[\rho_{q}, \rho_{q'}] = 0$ exactly on every window. So on $\mathcal{D}$, $[\rho_{q}, \rho_{q'}^{*}] = q\,\delta_{qq'}$ and $[\rho_{q}, \rho_{q'}] = 0$. These are the relations of Theorem \claimref{singularity.bosonisation-bridge-s-wave-sector}(1) on the infinite filled-sea array (there for the s-wave boundary modes); what is added is that the band modes of the null vacuum are such an array, and the explicit invariant domain $\mathcal{D}$. The check verifies, on $6$ and $8$ modes with the lower half filled, that $\rho_{1}, \rho_{2}$ annihilate the sea, that $[\rho_{q}, \rho_{q}^{*}]$ acts as $q$ on the sea, on $\rho_{1}^{*}$ of the sea ($q = 1, 2$) and on $(\rho_{1}^{*})^{2}$ of the sea ($q = 1$), and that $[\rho_{1}, \rho_{2}^{*}]$ and $[\rho_{1}, \rho_{2}]$ annihilate the depth-two states on $8$ modes. Finally, no finite window carries a scalar central term: by (v) the commutator is an edge operator of scalar part $0$, as Proposition \claimref{field-theory.no-canonical-commutation-relations-tower} requires of every floor.

\emph{(vii)} Here $\omega$ is a state of $\mathcal{A}_{e}$ and $\omega^{\otimes N}$ the product state on the even algebras of the $N$ blocks, which by (a) commute, so it is a state of the algebra they generate (for a block state that is not even, the product over the full block algebras need not be a state, since odd elements of distinct blocks anticommute). Constants commute with everything, so $[F_{N}(Q), F_{N}(R)] = N^{-1}\sum_{k,l}[Q_{k}, R_{l}]$, and by (a) only $k = l$ survives: $N^{-1}\sum_{k}C_{k}$. Put $c = \omega(C)$, $c' = \omega(C^{*})$ and $D_{N} = N^{-1}\sum_{k}(C_{k} - c)$. By the preliminaries $\omega^{\otimes N}(C_{k}) = c$, so $\omega^{\otimes N}(D_{N}) = 0$, and for $k \ne l$, by (b), $\omega^{\otimes N}((C_{k}^{*} - c')(C_{l} - c)) = (c' - c')(c - c) = 0$, while for $k = l$ it is $\omega(C^{*}C) - c'c$. Summing, $\omega^{\otimes N}(D_{N}^{*}D_{N}) = N^{-2}\cdot N(\omega(C^{*}C) - c'c)$, and $\omega^{\otimes N}(D_{N}^{*})\omega^{\otimes N}(D_{N}) = 0$, which is the stated variance (here $c' = c$, since $\omega(X^{*}) = \omega(X)$). It tends to $0$ as $N \to \infty$, which is the mean-square convergence of the commutator to $\sigma_{\omega}(Q, R)$. The count of canonical pairs is (iv). For a real state and self-adjoint $X, Y \in \mathcal{A}_{e}$, $[X, Y]^{*} = [Y, X] = -[X, Y]$ and $\omega(Z^{*}) = \omega(Z)$, so $\omega([X, Y]) = 0$: the form lives on pairs $(a, a^{*})$, e.g. $\sigma_{\omega}(b, b^{*}) = p_{1} - p_{2}$, and no unit is needed. A Weyl form with self-adjoint fields would need a central element of square $-1$ to turn $[X, Y]$ into a self-adjoint scalar, and the centre of $\mathcal{A}_{e}$ is spanned by $1$ and the block's pseudoscalar, of square $+1$. For a gauge-invariant product state $\omega(b) = 0$, so $F_{N}(b) = B$. The check verifies the identity, the mean and the variance exactly for $N = 1, 2, 3$, four pairs $(Q, R)$ (among them $b + b^{*}$ with $b - b^{*}$, $c + c^{*}$ with $n_{1}$, and pairs of off-diagonal matrix units) and three block states, with the single-block variance nonnegative and nonzero in most cases. The further statement of the quantum central limit theorem, that the fluctuation states themselves converge to a quasi-free state of the algebra of canonical commutation relations with this form, is not used.
